
\chapter{Conclusion}


\noindent The deployment of the Digital Admission Procedure Management system marks the successful transition of the application from development to a real-world operational environment. Through proper configuration of the backend and frontend components, the system has been effectively hosted on the university infrastructure, ensuring accessibility for end-users.

\vspace{0.3cm}

\noindent This deployment demonstrates the practicality and reliability of the system in managing the admission process in a more efficient and automated manner. It reduces manual workload, improves data handling accuracy, and ensures smoother communication between applicants and the administrative office.

\vspace{0.3cm}

\noindent Although the system is currently fully functional, there is still room for future enhancements such as mobile application support, improved scalability, and integration of advanced technologies like artificial intelligence. These improvements can further increase usability and performance.

\vspace{0.3cm}

\noindent Overall, the deployment phase confirms that the Digital Admission Procedure Management system is a successful step toward digitalizing the admission process at Jashore University of Science and Technology.

\vspace{0.3cm}

\noindent
Another key aspect brought out by the implementation process is the role of cooperation between the development team and the administrative body of the institution. It ensures that the software application fulfills all institutional needs and reflects the actual admission process.On another note, the software application has been created to be scalable and thus capable of being expanded without any need to redesign.